\documentclass[border=8pt]{standalone}
\usepackage{ctex}
\usepackage{amsmath}
\usepackage{tikz}
\begin{document}
\definecolor{mmRootFill}{RGB}{18,85,126}
\definecolor{mmLOneFill}{RGB}{234,148,136}
\begin{tikzpicture}[
  mmroot/.style={draw=black!75, line width=1.0pt, rounded corners=3.5pt, fill=mmRootFill, inner sep=5pt, font=\normalsize\bfseries},
  mml1/.style={draw=black!55, line width=0.9pt, rounded corners=2.5pt, fill=mmLOneFill, inner sep=4.5pt, font=\small\bfseries, align=left, text width=4.8cm},
  mml2/.style={draw=black!45, line width=0.8pt, rounded corners=2pt, fill=white, inner sep=4pt, font=\small, align=left, text width=6.2cm},
]
  %% 连线：根 → 一级 → 二级（节点填充不透明，盖住线头）
  \draw[black!55, line width=0.9pt] (4.35, 6.9) -- (4.8, 14.375);
  \draw[black!55, line width=0.9pt] (4.35, 6.9) -- (4.8, 13.225);
  \draw[black!55, line width=0.9pt] (4.35, 6.9) -- (4.8, 12.075);
  \draw[black!55, line width=0.9pt] (4.35, 6.9) -- (4.8, 10.925);
  \draw[black!55, line width=0.9pt] (4.35, 6.9) -- (4.8, 9.775);
  \draw[black!55, line width=0.9pt] (4.35, 6.9) -- (4.8, 8.625);
  \draw[black!55, line width=0.9pt] (4.35, 6.9) -- (4.8, 7.475);
  \draw[black!55, line width=0.9pt] (4.35, 6.9) -- (4.8, 6.325);
  \draw[black!55, line width=0.9pt] (4.35, 6.9) -- (4.8, 5.175);
  \draw[black!55, line width=0.9pt] (4.35, 6.9) -- (4.8, 4.025);
  \draw[black!55, line width=0.9pt] (4.35, 6.9) -- (4.8, 2.875);
  \draw[black!55, line width=0.9pt] (4.35, 6.9) -- (4.8, 1.725);
  \draw[black!55, line width=0.9pt] (4.35, 6.9) -- (4.8, 0.575);
  \draw[black!55, line width=0.9pt] (4.35, 6.9) -- (4.8, -0.575);
  %% 节点
  \node[mmroot, anchor=east] at (4.35, 6.9) {2.8 直线与圆锥曲线 \\ 的位置关系};
  \node[mml1, anchor=west] at (4.6, 14.375) {直线与椭圆位置关系：联立消去y得一个关于x的一元二次方程．};
  \node[mml1, anchor=west] at (4.6, 13.225) {圆锥曲线的弦与弦长};
  \node[mml1, anchor=west] at (4.6, 12.075) {直线与圆锥曲线位置关系的判定与相切的辨析};
  \node[mml1, anchor=west] at (4.6, 10.925) {圆锥曲线的中点弦斜率结论};
  \node[mml1, anchor=west] at (4.6, 9.775) {硬解定理};
  \node[mml1, anchor=west] at (4.6, 8.625) {第二焦半径公式};
  \node[mml1, anchor=west] at (4.6, 7.475) {焦点弦定理};
  \node[mml1, anchor=west] at (4.6, 6.325) {阿基米德三角形};
  \node[mml1, anchor=west] at (4.6, 5.175) {蒙日圆};
  \node[mml1, anchor=west] at (4.6, 4.025) {定比分点};
  \node[mml1, anchor=west] at (4.6, 2.875) {定比点差定理};
  \node[mml1, anchor=west] at (4.6, 1.725) {齐次式};
  \node[mml1, anchor=west] at (4.6, 0.575) {齐次方程};
  \node[mml1, anchor=west] at (4.6, -0.575) {非对称问题};
\end{tikzpicture}
\end{document}
